% from ega4-8-fr.tex
\oldpage[IV-3]{20}
Nous poserons alors
\begin{equation}
  \mathcal F=u_\lambda^*(\mathcal F_\lambda)
  \tag{8.5.1.2}\label{IV.8.5.1.2-fr}
\end{equation}
qui est un $\mathcal O_S$-Module ne dépendant pas de l'indice $\lambda\in L$,
en vertu de l'hypothèse (8.5.1.1).

Soient maintenant $(\mathcal F_\lambda)$, $(\mathcal G_\lambda)$ deux telles
familles de $\mathcal O_{S_\lambda}$-Modules. Il est clair que les applications
$f_\lambda\mapsto u_{\lambda\mu}^*(f_\lambda)$ de
$\operatorname{Hom}_{S_\lambda}(\mathcal F_\lambda,\mathcal G_\lambda)$ dans
$\operatorname{Hom}_{S_\mu}(\mathcal F_\mu,\mathcal G_\mu)$ définissent alors
un \emph{système inductif} de groupes abéliens
$(\operatorname{Hom}_{S_\lambda}(\mathcal F_\lambda,\mathcal G_\lambda),
u_{\lambda\mu}^*)$, et que les applications
$f_\lambda\mapsto u_\lambda^*(f_\lambda)$ forment un \emph{système inductif}
d'homomorphismes de groupes abéliens
\[
  u_\lambda^*:\operatorname{Hom}_{S_\lambda}
  (\mathcal F_\lambda,\mathcal G_\lambda)
  \longrightarrow\operatorname{Hom}_S(\mathcal F,\mathcal G)~;
\]
d'où, en passant à la limite inductive, un homomorphisme canonique de groupes
abéliens
\begin{equation}
  u_{\mathcal F,\mathcal G}:
  \varinjlim\operatorname{Hom}_{S_\lambda}
  (\mathcal F_\lambda,\mathcal G_\lambda)
  \longrightarrow\operatorname{Hom}_S(\mathcal F,\mathcal G).
  \tag{8.5.1.3}\label{IV.8.5.1.3-fr}
\end{equation}

Notons que lorsque $\mathcal F_\lambda=\mathcal O_{S_\lambda}$, la condition
(8.5.1.1) est vérifiée, et que l'on a $\mathcal F=\mathcal O_S$~;
l'homomorphisme (8.5.1.3) donne donc $(0_{\mathrm I},5.1.1)$ un homomorphisme
canonique de groupes abéliens
\begin{equation}
  u_{\mathcal G}:\varinjlim\Gamma(S_\lambda,\mathcal G_\lambda)
  \longrightarrow\Gamma(S,\mathcal G).
  \tag{8.5.1.4}\label{IV.8.5.1.4-fr}
\end{equation}
\end{env}

\begin{theorem}[8.5.2]
\label{IV.8.5.2-fr}
\begin{enumerate}[label=(\roman*)]
  \item Supposons $S_0$ quasi-compact (resp. quasi-compact et quasi-séparé) et
  que, pour un $\lambda\in L$, $\mathcal F_\lambda$ soit quasi-cohérent et de
  type fini (resp. de présentation finie) et $\mathcal G_\lambda$
  quasi-cohérent. Alors l'homomorphisme $u_{\mathcal F,\mathcal G}$ est
  injectif (resp. bijectif).

  \item Supposons $S_0$ quasi-compact et quasi-séparé. Pour tout
  $\mathcal O_S$-Module quasi-cohérent $\mathcal F$ de présentation finie, il
  existe un $\lambda\in L$ et un $\mathcal O_{S_\lambda}$-Module
  quasi-cohérent $\mathcal F_\lambda$ de présentation finie tels que
  $\mathcal F$ soit isomorphe à $u_\lambda^*(\mathcal F_\lambda)$.
\end{enumerate}
\end{theorem}

\begin{proof}
\begin{enumerate}[label=(\roman*)]
  \item On peut évidemment se borner au cas où $S_0=S_\lambda$ puisque les
  morphismes $u_{0\lambda}:S_\lambda\to S_0$ sont affines, donc quasi-compacts
  et séparés. Considérons d'abord le cas où
  $S_0=\operatorname{Spec}(A_0)$ est \emph{affine}. Alors l'assertion (i)
  équivaut au

  \begin{lemma}[8.5.2.1]
  \label{IV.8.5.2.1-fr}
  Soient $A_0$ un anneau, $(A_\lambda)_{\lambda\in L}$ un système inductif de
  $A_0$-algèbres, $A=\varinjlim A_\lambda$~; soient $M_0$, $N_0$ deux
  $A_0$-modules, et posons
  $M_\lambda=M_0\otimes_{A_0}A_\lambda$,
  $N_\lambda=N_0\otimes_{A_0}A_\lambda$,
  $M=M_0\otimes_{A_0}A=\varinjlim M_\lambda$,
  $N=N_0\otimes_{A_0}A=\varinjlim N_\lambda$. Si $M_0$ est de type fini
  (resp. de présentation finie), l'homomorphisme canonique
  \begin{equation}
    \varinjlim\operatorname{Hom}_{A_\lambda}(M_\lambda,N_\lambda)
    \longrightarrow\operatorname{Hom}_A(M,N)
    \tag{8.5.2.2}\label{IV.8.5.2.2-fr}
  \end{equation}
  est injectif (resp. bijectif).
  \end{lemma}

  On sait en effet (Bourbaki, \emph{Alg.}, chap. II, 3\textsuperscript{e} éd.,
  \S~5, n\textsuperscript{o}~1) que l'on a des isomorphismes canoniques
  fonctoriels
  \[
    \operatorname{Hom}_{A_\lambda}(M_\lambda,N_\lambda)
    \xrightarrow{\sim}\operatorname{Hom}_{A_0}(M_0,N_\lambda),
    \qquad
    \operatorname{Hom}_A(M,N)
    \xrightarrow{\sim}\operatorname{Hom}_{A_0}(M_0,N)
  \]

  